\phantomsection\addcontentsline{toc}{section}{Robotics}
\ECchapter{07}{Robotics}{How can robots work safely with humans?}{ECBlue}

\ECObjectives{
\begin{itemize}
\item Identify the mechanical, electrical, control and safety subsystems of a robot.
\item Explain position, force and vision feedback in a collaborative workstation.
\item Design and justify a safe robotic solution for an industrial task.
\end{itemize}}

\begin{multicols}{2}
\begin{engineeringnews}
Collaborative robots, or \textbf{cobots}, are changing the organisation of factories and laboratories. Instead of operating permanently behind a fence, they can share selected tasks with people. The engineering objective is not merely to make the robot slower: the complete workstation must detect hazards, limit energy and bring the system to a safe state when required.
\end{engineeringnews}

\begin{ecarticle}
A robot is a programmable machine that senses its environment, processes information and acts through a mechanical structure. A typical articulated robot contains rigid links connected by joints. Electric servomotors create torque at these joints, while encoders measure angular position and speed. A controller compares the measured motion with the requested trajectory and continuously corrects the motor commands.

The tool attached to the final link is called the \textbf{end-effector}. It may be a gripper, a welding torch, a screwdriver or a camera. Selecting this tool requires engineers to consider payload, accuracy, cycle time, gripping force and product geometry. The useful region that the end-effector can reach is the robot's workspace.

Position control is sufficient for many repetitive tasks, but physical interaction requires additional information. A force--torque sensor can detect contact, while motor-current monitoring can estimate resistance at a joint. Machine-vision systems locate parts, inspect surfaces and track people. These measurements form feedback loops that allow the robot to adapt rather than repeat a completely fixed sequence.

Safe collaboration depends on the entire application. Possible strategies include monitored stops, hand-guiding, speed and separation monitoring, and power-and-force limitation. A scanner or camera can define protective zones around the robot. As a person approaches, the controller first reduces speed and then stops hazardous motion. Mechanical design also matters: rounded surfaces, low moving mass and the avoidance of crushing points can reduce injury risk.

Mobile robots combine localisation, mapping and path planning. Because sensors can fail or be obstructed, their design requires diagnostics and a defined safe response. Cybersecurity is equally important: access control, validated software and network segmentation help prevent dangerous commands or corrupted sensor data.
\end{ecarticle}

\begin{tcolorbox}[title=\sffamily\bfseries Engineering Insight,colback=yellow!10,colframe=ECOrange]
A cobot is not automatically a safe collaborative application. Safety depends on the robot, tool, workpiece, programmed speed, human task and surrounding equipment. The risk assessment must therefore cover the complete workstation.
\end{tcolorbox}

\columnbreak

\begin{engineeringfocus}
\textbf{Closed-loop architecture of a collaborative robot}

\begin{center}
\resizebox{\linewidth}{!}{%
\begin{tikzpicture}[node distance=5mm and 5mm,font=\scriptsize,>=Latex]
\tikzset{block/.style={draw=ECBlue,fill=ECLightBlue,rounded corners,align=center,minimum height=9mm,minimum width=18mm}}
\node[block] (plan) {task and\\trajectory};
\node[block,right=of plan] (ctrl) {robot\\controller};
\node[block,right=of ctrl] (drive) {servo\\drives};
\node[block,right=of drive] (mech) {joints and\\end-effector};
\node[block,below=8mm of drive] (sense) {encoders, force,\\vision, scanner};
\node[draw=ECOrange,fill=ECOrange!12,rounded corners,below=8mm of plan,align=center,minimum height=9mm] (safe) {safety\\controller};
\draw[->,very thick,ECBlue] (plan) -- (ctrl);
\draw[->,very thick,ECBlue] (ctrl) -- (drive);
\draw[->,very thick,ECBlue] (drive) -- (mech);
\draw[->,very thick,ECGreen] (mech.south) |- (sense.east);
\draw[->,very thick,ECGreen] (sense.west) -| (ctrl.south);
\draw[->,very thick,ECOrange] (sense) -- (safe);
\draw[->,very thick,ECOrange] (safe) -- (ctrl);
\end{tikzpicture}%
}
\end{center}

The standard motion controller aims to achieve the task. The safety controller independently supervises hazardous conditions and can request a controlled stop or remove drive power.
\end{engineeringfocus}

\begin{keyfigures}
\begin{tabularx}{\linewidth}{@{}Xr@{}}
Typical articulated cobot axes & 6--7\\
Repeatability of small industrial robots & about 0.02--0.1 mm\\
Typical cobot payload range & 3--35 kg\\
Main feedback period & milliseconds\\
Emergency response objective & safe state
\end{tabularx}
\end{keyfigures}

\begin{tcolorbox}[title=\sffamily\bfseries Assembly cycle comparison,colback=white,colframe=ECBlue]
\begin{center}
\begin{tikzpicture}
\begin{axis}[width=.96\linewidth,height=45mm,ybar,bar width=13pt,ymin=0,ymax=60,ylabel={Cycle time (s)},symbolic x coords={Manual,Cobot-assisted,Industrial robot},xtick=data,xticklabel style={font=\scriptsize,align=center,rotate=12,anchor=east},yticklabel style={font=\scriptsize},grid=major,nodes near coords,every node near coord/.append style={font=\scriptsize}]
\addplot[fill=ECBlue] table[x=mode,y=cycle_time,col sep=comma]{data/EC07/robot_cycle.csv};
\end{axis}
\end{tikzpicture}
\end{center}
\footnotesize Illustrative values for one assembly operation. A faster cycle does not by itself prove that a solution is safer, more flexible or more economical.
\end{tcolorbox}

\begin{vocabulary}
\begin{tabularx}{\linewidth}{@{}lX@{}}
joint & articulation\\
link & solide / bras\\
actuator & actionneur\\
encoder & codeur\\
end-effector & effecteur terminal\\
payload & charge utile\\
repeatability & répétabilité\\
protective stop & arrêt de protection
\end{tabularx}
\end{vocabulary}

\begin{questions}
\item What information is provided by an encoder?
\item Distinguish accuracy from repeatability.
\item Why may position control be insufficient during contact tasks?
\item Explain speed and separation monitoring.
\item Why must the end-effector be included in the risk assessment?
\item Give two possible consequences of a cybersecurity failure.
\item Use the graph to explain why productivity cannot be the only selection criterion.
\end{questions}

\begin{engineeringchallenge}
A company wants to automate the insertion and tightening of four screws on a 2.5~kg product while an operator performs the adjacent wiring task. Design the collaborative workstation. Select the robot, end-effector and sensors; define the work sequence and protective zones; identify crushing, impact, tool and cybersecurity hazards; and propose validation tests. Compare the solution with a fully manual station and a fenced industrial robot using a weighted decision matrix.
\end{engineeringchallenge}

\begin{applicationexercise}
A collaborative robot moves a 4.0~kg component with a horizontal acceleration of 1.8~m/s$^2$. Calculate the resulting inertial force. If the end-effector has a mass of 1.5~kg, repeat the calculation for the complete moving load. Explain why this simple calculation is not sufficient to validate human--robot safety.
\end{applicationexercise}

\begin{speaking}
Prepare a four-minute design review answering the question: ``Should this workstation use a cobot or a conventional industrial robot?'' Defend your choice using safety, cycle time, flexibility, payload, maintenance and cost.
\end{speaking}
\end{multicols}

\subsection*{Sources}
\ECsource{International Federation of Robotics -- industrial and service robots}{https://ifr.org/}
\ECsource{NIST -- robotics and autonomous systems}{https://www.nist.gov/topics/robotics}
\ECsource{European Commission -- machinery safety}{https://single-market-economy.ec.europa.eu/sectors/mechanical-engineering/machinery_en}
